\chapter*{Lời cảm ơn}
\label{thanks}

\begin{spacing}{1.15}
Lời đầu tiên, chúng em xin gửi lời cảm ơn chân thành và sâu sắc nhất đến Trường Đại học Khoa học Tự nhiên, ĐHQG-HCM, đặc biệt là Khoa Công nghệ Thông tin đã tạo điều kiện học tập và môi trường nghiên cứu tốt nhất cho chúng em trong suốt thời gian qua.

Đặc biệt, chúng em xin gửi lời tri ân sâu sắc đến \textbf{ThS. Phạm Trọng Nghĩa} và \textbf{ThS. Lê Nhựt Nam} - hai thầy hướng dẫn đã luôn tận tình chỉ bảo, định hướng và cung cấp những tài liệu, kiến thức quý báu giúp chúng em hoàn thành khóa luận này. Sự kiên nhẫn và tâm huyết của các thầy là nguồn động lực to lớn giúp chúng em vượt qua những khó khăn trong quá trình nghiên cứu về Học máy Lượng tử.

Bên cạnh đó, chúng em cũng xin cảm ơn các thầy cô trong hội đồng đánh giá đã dành thời gian đọc và góp ý cho cuốn khóa luận này để chúng em nhận ra những thiếu sót và tiếp tục hoàn thiện hơn trong tương lai.

Cuối cùng, chúng em xin gửi lời biết ơn sâu sắc đến gia đình, người thân và bạn bè - những người đã luôn ở bên cạnh động viên, cổ vũ và tạo mọi điều kiện tốt nhất về cả vật chất lẫn tinh thần để chúng em yên tâm học tập và nghiên cứu.

Dù đã rất cố gắng hoàn thiện, nhưng do hạn chế về thời gian cũng như kiến thức và kinh nghiệm thực tế, cuốn khóa luận này khó tránh khỏi những thiếu sót. Chúng em rất mong nhận được sự lượng thứ và những đóng góp quý báu từ các thầy cô và bạn đọc.

\vspace{0.4cm}
\begin{flushright}
\textit{Thành phố Hồ Chí Minh, tháng 07/2026}\\
\textbf{Nhóm sinh viên thực hiện}\\
Nguyễn Sinh Trực - Phạm Tuấn Vương
\end{flushright}
\end{spacing}